\chapter*{Anhang: Verwendete KI-basierte elektronische Hilfsmittel}
\vspace{-6em}{\Huge\color{dkfzgray2} Appendix: AI-based tools used}
\addcontentsline{toc}{chapter}{Anhang: Verwendete KI-basierte elektronische Hilfsmittel}

Im Folgenden werden die im Rahmen dieser Dissertation verwendeten KI-basierten
Hilfsmittel gemäß den Vorgaben aus dem Anhang „Angabe zu verwendeten KI-basierten
elektronischen Hilfsmitteln" dokumentiert.

\emph{\small\color{dkfzgray2} The AI-based tools used in the course of this
dissertation are documented below, in accordance with the requirements set out
in the appendix ``Information on the use of AI-based tools''.}

\bigskip

% ---------------------------------------------------------------------------
% ANLEITUNG (vor Einreichung löschen):
%
% Für JEDES verwendete KI-Tool (z. B. Claude / Claude Code, ChatGPT, DeepL,
% Grammarly, GitHub Copilot, ...) einen eigenen Block nach untenstehendem
% Muster anlegen. Pflichtangaben laut KI-Leitlinie der Universität Heidelberg:
%   1. Ziel:              Wofür wurde das Tool eingesetzt?
%   2. Verwendungsweise:  Wie genau wurde es eingesetzt? (z. B. Grammatik-
%                         korrektur, Übersetzung, Textentwurf, Code-Erstellung,
%                         Literaturrecherche) UND die Bestätigung, dass jeder
%                         Output eigenständig fachlich geprüft wurde.
%   3. Betroffene Kapitel/Abschnitte: konkret benennen, nicht "überall".
%
% Wurde gar keine KI verwendet, hier explizit "Es wurden keine KI-basierten
% Hilfsmittel verwendet." eintragen statt das Kapitel leer zu lassen.
% ---------------------------------------------------------------------------

\section*{TODO: Name des Tools (z.\,B. Claude Code)}
\addcontentsline{toc}{section}{TODO: Name des Tools}

\begin{enumerate}
    \item \textbf{Ziel.} TODO: Wofür wurde das Tool eingesetzt (z.\,B.
    Rechtschreib-/Grammatikprüfung, Literaturrecherche, Umformulierung
    bestehenden Textes, Erstellung von Textentwürfen, Code-Generierung für
    Abbildungen, Konvertierung eines bestehenden Entwurfs in dieses
    LaTeX-Template)?\\
    \emph{\small\color{dkfzgray2} Goal. TODO: What was the tool used for?}

    \item \textbf{Verwendungsweise.} TODO: Wie wurde das Tool konkret
    eingesetzt? Wichtig: hier auch festhalten, dass alle Vorschläge/Entwürfe
    vom Verfasser/von der Verfasserin eigenständig fachlich geprüft, korrigiert
    und erst danach übernommen wurden.\\
    \emph{\small\color{dkfzgray2} How used. TODO: How exactly was the tool
    used? State that all suggestions/drafts were independently reviewed,
    corrected, and approved by the author before adoption.}

    \item \textbf{Betroffene Kapitel/Abschnitte.} TODO: Konkrete Kapitel/
    Abschnitte benennen (z.\,B. "Einleitung, Abschnitt 2.3, Abbildungen 3.4
    und 3.7").\\
    \emph{\small\color{dkfzgray2} Chapters/sections affected. TODO: Name the
    specific chapters/sections.}
\end{enumerate}

\bigskip
\textit{TODO: diesen Block für jedes weitere verwendete KI-Tool duplizieren.}
